% PAGES: 139-140
% Continues sec04_c2_2.tex (pages 136-138), whose final \end{proof} (proof
% of Theorem 4.4) closes cleanly at the bottom of page 138 -- confirmed by
% page 139 opening with plain narrative text ("The result below is an
% equivalent characterization...") rather than a "Proof." continuation.
% No environment is left open across that boundary.
%
% Still section 4.2 "Diagonal forms" through Lemma 4.8; section 4.3
% "Diagonally dominant matrices" begins near the bottom of page 139 (folio
% 123) with the R_j definition (4.28).
%
% This is the last file of the c2 chunk (PDF pages 134-140 / folios
% 118-124). Theorem 4.6 (Gershgorin's Theorem) is stated in full at the
% end of page 140 (folio 124), followed by a complete explanatory remark
% ("Note that Gershgorin's Theorem can also be stated as follows: ...in
% the complex space."). That remark sentence is grammatically complete and
% is the last thing on page 140 -- no \[ \], \begin{...}, or sentence is
% left open. The theorem's "Proof." itself has not yet appeared on the
% page (it begins on the next page, folio 125, PDF page 141, outside this
% assignment); the next chunk's agent should open it with a plain
% \begin{proof}...\end{proof} following this file's content, not as a
% continuation of anything opened here.

The result below is an equivalent characterization of diagonalizable matrices and
follows immediately from the proof of Theorem 4.4 above:

\begin{theorem}
A square matrix of size $n$ has a diagonal form if and only if the matrix has $n$
linearly independent eigenvectors.
\end{theorem}

\begin{lemma}
Let $A$ be a diagonalizable matrix, and let $A = V\Lambda V^{-1}$ be its diagonal
form. For every positive integer $p$, the matrix $A^p$ is diagonalizable and has the
following diagonal form:
\begin{equation}
A^p \;=\; V\Lambda^p V^{-1}.
\end{equation}
\end{lemma}

\begin{proof}
If $p$ is a positive integer, then
\begin{equation}
\begin{split}
A^p &= (V\Lambda V^{-1})^p \;=\; (V\Lambda V^{-1})(V\Lambda V^{-1}) \dots (V\Lambda
V^{-1})(V\Lambda V^{-1})\\
&= V\Lambda (V^{-1}V) \Lambda (V^{-1}V) \dots (V^{-1}V) \Lambda V^{-1}\\
&= V\Lambda^p V^{-1},
\end{split}
\end{equation}
since $V^{-1}V = I$.
\end{proof}

\begin{lemma}
Let $A$ be a nonsingular diagonalizable matrix, and let $A = V\Lambda V^{-1}$ be its
diagonal form. Then, the inverse matrix $A^{-1}$ is also diagonalizable and has the
diagonal form $A^{-1} = V\Lambda^{-1}V^{-1}$. Moreover, for every positive integer
$p$, the matrix $A^{-p}$ is diagonalizable and has the following diagonal form:
\[
A^{-p} \;=\; V\Lambda^{-p} V^{-1}.
\]
\end{lemma}

\begin{proof}
From Lemma 1.7, it follows that
\[
A^{-1} \;=\; (V\Lambda V^{-1})^{-1} \;=\; (V^{-1})^{-1}\Lambda^{-1}V^{-1} \;=\;
V\Lambda^{-1}V^{-1}.
\]
By definition, if $p$ is a positive integer, then $A^{-p} = (A^p)^{-1}$. Using
(4.26), we obtain that
\begin{align*}
A^{-p} &= (A^p)^{-1} \;=\; (V\Lambda^p V^{-1})^{-1} \;=\; (V^{-1})^{-1}(\Lambda^p)^{-1}
V^{-1}\\
&= V\Lambda^{-p}V^{-1}.
\end{align*}
\end{proof}

\section{Diagonally dominant matrices}

Let $A$ be an $n \times n$ matrix, and denote by
\begin{equation}
R_j \;=\; \sum_{k=1:n,\, k \neq j} |A(j,k)|
\end{equation}
the sum of the absolute values of all the entries on the $j$-th row of the matrix $A$
except for the main diagonal entry $A(j,j)$.

\begin{definition}
The $n \times n$ matrix $A$ is weakly diagonally dominant if and only if
\begin{equation}
|A(j,j)| \;\geq\; R_j, \quad \forall\, j = 1:n,
\end{equation}
where $R_j$ is given by (4.28), i.e., if and only if, for every row, the absolute
value of the main diagonal entry from the row is greater than the sum of the absolute
values of all the other entries in that row.
\end{definition}

\begin{definition}
The $n \times n$ matrix $A$ is strictly diagonally dominant if and only if
\begin{equation}
|A(j,j)| \;>\; R_j, \quad \forall\, j = 1:n,
\end{equation}
where $R_j$ is given by (4.28).
\end{definition}

\begin{examplex}
The $N \times N$ matrix
\[
B_N \;=\; \begin{pmatrix}
2 & -1 & \dots & 0\\
-1 & \ddots & \ddots & \vdots\\
\vdots & \ddots & \ddots & -1\\
0 & \dots & -1 & 2
\end{pmatrix}
\]
is a weakly diagonally dominant matrix, since the only nonzero entries of $B_N$ are
\begin{align*}
B_N(j,j) &= 2, \quad \forall\, j = 1:N;\\
B_N(j,j-1) &= -1, \quad \forall\, j = 2:N;\\
B_N(j,j+1) &= -1, \quad \forall\, j = 1:(N-1),
\end{align*}
and therefore
\begin{align*}
\sum_{k=2:N} |B_N(1,k)| &= |B_N(1,2)| = 1 < 2 = |B_N(1,1)|;\\
\sum_{k=1:(N-1)} |B_N(N,k)| &= |B_N(N,N-1)| = 1 < 2 = |B_N(N,N)|;\\
\sum_{k=1:N,\, k \neq j} |B_N(j,k)| &= |B_N(j,j-1)| + |B_N(j,j+1)| = 2 = |B_N(j,j)|,
\end{align*}
% UNCLEAR: p.140 prints the range as $j = 2:(N-2)$; the interior rows of
% $B_N$ run $j = 2:(N-1)$. Author error, reproduced verbatim.
for all $j = 2:(N-2)$.

Then, from definition 4.7, it follows that $B_N$ is a weakly diagonally dominant
matrix.
\end{examplex}

\begin{theorem}
(Gershgorin's Theorem.) Let $A$ be a square matrix of size $n$. For any eigenvalue
$\lambda$ of $A$, there
exists an index $j$, with $1 \leq j \leq n$, such that
\begin{equation}
|\lambda - A(j,j)| \;\leq\; R_j,
\end{equation}
where $R_j$ is given by (4.28).
\end{theorem}

Note that Gershgorin's Theorem can also be stated as follows:
\[
\lambda(A) \;\subset\; \bigcup_{j=1}^n D\bigl(A(j,j), R_j\bigr),
\]
where $\lambda(A)$ is the set of all the eigenvalues of $A$ and $D(a,R) = \{z \in \C
\ \text{such that} \ |z-a| \leq R\}$ is the disc of center $a \in \C$ and radius
$R > 0$ in the complex space.
